% Generated draft: requires mathematical and source review.
\documentclass[11pt]{article}
\usepackage[margin=25mm]{geometry}
\usepackage{fontspec}
\setmainfont{Latin Modern Roman}
\usepackage{amsmath,amssymb,mathtools,unicode-math}
\setmathfont{Latin Modern Math}
\usepackage{enumitem,needspace,xurl,hyperref}
\setlength{\parindent}{0pt}
\setlength{\parskip}{5pt}
\setlength{\emergencystretch}{3em}
\newcommand{\N}{\mathbb{N}}
\newcommand{\Z}{\mathbb{Z}}
\newcommand{\Q}{\mathbb{Q}}
\newcommand{\R}{\mathbb{R}}
\newcommand{\C}{\mathbb{C}}
\newcommand{\F}{\mathbb{F}}
\newcommand{\A}{\mathbb{A}}
\newcommand{\PP}{\mathbb{P}}
\newcommand{\E}{\mathbb{E}}
\newcommand{\eps}{\varepsilon}
\newcommand{\dd}{\,\mathrm{d}}
\begin{document}
\section*{UnsolvedMath ID: 700945}
\textbf{Source problem code:} catalogue entry
\subsection{Definitions and mathematical statement}
For subgroups \(A,B\leq G\), write \(AB=\{ab:a\in A,\ b\in B\}\). For a prime \(p\), the quasicyclic \(p\)-group is the additive group \(C_{p^\infty}=\Z[1/p]/\Z\). A group is called a Chernikov group if it has a normal subgroup of finite index isomorphic to \(\prod_{i=1}^{r}C_{p_i^\infty}\), where \(r\geq 0\) is finite and the \(p_i\) are primes, with repetitions allowed. The empty product is the trivial group, so finite groups are included. A group \(H\) is polycyclic if it has a finite subnormal series \(1=H_0\triangleleft H_1\triangleleft\cdots\triangleleft H_n=H\) whose factors \(H_i/H_{i-1}\) are cyclic, possibly finite. Here ``almost polycyclic'' is interpreted in its standard sense: having a polycyclic subgroup of finite index.

For every group \(G\) and every three subgroups \(A,B,C\leq G\) satisfying the equalities of sets
\[
G=AB=AC=BC,
\]
ask the following two questions separately:
\begin{enumerate}
\item[(a)] If \(A,B,C\) are Chernikov groups, must \(G\) be a Chernikov group?
\item[(b)] If \(A,B,C\) are almost polycyclic, must \(G\) be almost polycyclic?
\end{enumerate}
No finiteness, finite generation, normality, or solvability assumption on \(G\) or on these subgroups is supplied beyond the hypotheses in each question. The source does not explicitly define the two group classes; the conventional definitions above specify the interpretation used in this draft. Attribution remains unresolved: the statement begins with ``(B. Amberg)'', whereas the background and proposer field name L. S. Kazarin.
\subsection{Short English statement}
If a group is the product of each pair among three subgroups, must it be Chernikov when all three subgroups are Chernikov, and must it be almost polycyclic when all three are almost polycyclic? These are separate questions. The supplied registry labels the problem open; that status has not been independently verified.
\subsection{Sources}
\begin{itemize}
\item Local UnsolvedMath registry, supplied complete record with permanent ID \(700945\), exact source code \texttt{catalogue entry}, title ``Kourovka Notebook Problem 13.27''. The supplied record labels the problem open and records an update date of \texttt{2026-10-16}.
\item Supplied bibliographic attribution: \emph{The Kourovka Notebook (Unsolved Problems in Group Theory)}, problem 13.27, originating in the 13th issue (1995), as included in the cited 21st issue edition (2026), arXiv:1401.0300v48, \url{https://arxiv.org/abs/1401.0300v48}. The supplied collection description names E. I. Khukhro and V. D. Mazurov as editors. This reference has not been independently consulted.
\item The supplied background and proposer field attribute the problem to L. S. Kazarin and give the proposal year as 1995. The supplied statement instead carries the parenthetical attribution ``(B. Amberg)''; the record does not explain their respective roles.
\end{itemize}
\end{document}
